\section{Verification Snapshot}
The submission snapshot was checked on Windows using Node.js 24.19 and the
bundled Python environment on 17 September 2026. The web production build and
TypeScript compilation succeeded. The automated suites passed 41 AI tests,
9 web-domain tests, and 11 backend authorization/validation tests. Backend tests
use an isolated database stub; no live PostgreSQL or cloud deployment is claimed.

\begin{longtable}{p{0.20\textwidth}p{0.43\textwidth}p{0.27\textwidth}}
\toprule\textbf{Area} & \textbf{Verification evidence} & \textbf{Result / limitation}\\\midrule\endhead
AI & Recommendation constraints, retrieval, refusal on insufficient evidence, scan quality and HTTP contracts & 41 tests passed; local and synthetic datasets.\\
Web domain & Order transitions, service validation, scan boundaries, publication, approval/moderation, CSV safety and saved-state round trip & 9 tests passed.\\
Business API & Shared JWT key, malformed/expired/forged tokens, role checks, private order listing, safe login and server-derived order prices & 11 tests passed; database stub.\\
React build & TypeScript and Vite production output & Passed.\\
Mobile & Source and navigation/widget test inspected; booking estimate and demo wording corrected & Flutter SDK unavailable locally; device build not verified.\\
Report & LaTeX sources consolidated and compiled using Tectonic & PDF generated; visual review required for release.\\
\bottomrule
\end{longtable}

\section{Requirement-to-Code Examples}
\begin{description}[style=nextline]
\item[US-03 / SCRUM-29: inspect an order] A lab operator can search by order,
customer or service and filter by status. \path{apps/web/src/App.tsx} contains
the order table, selection and detail view. Keyboard Enter selects a row.
\item[US-04: advance a workflow] \path{apps/web/src/domain.ts} permits only
the next stage. Completed orders cannot advance, and scan orders at quality
check require reviewed publication before completion.
\item[US-05: deliver scans] JPEG/TIFF extension and MIME checks, positive size,
50 MB per-file limit, duplicate filenames and a 100-file limit are enforced.
The demo saves private metadata after quality and release confirmation;
actual file bytes remain on the device.
\item[US-08 and US-09 / SCRUM-30: administration] Incomplete lab applications
cannot be approved. Rejections and moderation decisions require reasons.
Repeated final decisions are rejected and local audit entries record changes.
Browser audit data is not tamper-proof server evidence.
\item[Security boundary] The Express API shares one configured signing and
verification secret, restricts user/schema inspection to administrators, and
filters non-admin orders by owner. Order totals come from the service catalogue.
\end{description}

\section{Workflow and Sequence Diagrams}
The diagrams below distinguish implemented local interactions from the target
integration design. The complete story/rule matrix and manual checklist are in
\path{docs/AcceptanceTraceability.md} and \path{docs/SubmissionChecklist.md}.

\begin{figure}[htbp]\centering
\begin{tikzpicture}[node distance=6mm, box/.style={draw,rounded corners,fill=teal!6,text width=11cm,align=center,minimum height=8mm,font=\small}, >=Stealth]
\node[box] (a) {New order: inspect customer, service and quantity};
\node[box,below=of a] (b) {Film received $\rightarrow$ Developing $\rightarrow$ Scanning};
\node[box,below=of b] (c) {Quality check: select valid JPEG/TIFF files};
\node[box,below=of c] (d) {Confirm quality review and release for the selected order};
\node[box,below=of d] (e) {Save private metadata + complete order + append local audit event};
\draw[->] (a)--(b);\draw[->] (b)--(c);\draw[->] (c)--(d);\draw[->] (d)--(e);
\end{tikzpicture}
\caption{Implemented local web workflow; no cloud upload or customer message is sent}
\end{figure}

\begin{figure}[htbp]\centering
\begin{tikzpicture}[x=1cm,y=1cm,>=Stealth,font=\small]
\node[draw,fill=teal!6,minimum width=2.8cm] at (0,0) {Browser form};
\node[draw,fill=teal!6,minimum width=2.8cm] at (5,0) {Domain rules};
\node[draw,fill=teal!6,minimum width=2.8cm] at (10,0) {Browser storage};
\foreach \x in {0,5,10}{\draw[dashed,gray] (\x,-0.3)--(\x,-4.7);}
\draw[->] (0,-0.9)--node[above]{Submit service / scan decision}(5,-0.9);
\draw[->] (5,-1.6)--node[above]{Validate; compute next state}(10,-1.6);
\draw[->,dashed] (10,-2.3)--node[above]{Saved JSON or storage error}(5,-2.3);
\draw[->,dashed] (5,-3.0)--node[above]{Update UI only on successful save}(0,-3.0);
\draw[->] (0,-3.8)--node[above]{Reload: read and validate snapshot}(10,-3.8);
\end{tikzpicture}
\caption{Implemented state persistence, validation and error handling in the web demo}
\end{figure}

\clearpage
\section{Known Limits and Demonstration Talking Points}
The assessed implementation is a documented prototype with executable local
interactions and testable backend/AI modules. It is not a completed production
marketplace. Shared frontend authentication, live database migration verification,
cross-device synchronization, cloud scan storage, payment/logistics adapters,
mobile device testing and deployment remain explicit integration tasks.

For a five-minute demonstration: introduce the problem and actors; show order
search and a valid status transition; demonstrate a service edit that survives
reload; show scan validation and private publication metadata; demonstrate an
admin decision and its audit event; finish with the test summary and the
requirement-to-code matrix. Explain that local role switching is a demonstration
control and that backend authorization must independently enforce permissions.
